% ============================================================
%  表 III（中文伴生版）：与 ../tables/tab-p2-time.tex 数值 100% 同步
%  数据源：python/experiments/p5_freeze_20261124/p5_p2_h5.csv
% ============================================================
\begin{table}[t]
\centering
\caption{固定次信道时间参考留出残差与单侧馈线断开恢复实例（轮次 2715--2717，\texttt{time} 配置）}
\label{tab:p2_time}
\setlength{\tabcolsep}{1.2pt}
\scriptsize
\begin{tabular}{@{}llccccccc@{}}
\toprule
\textbf{轮} & \textbf{场景} & \textbf{联合} & \multicolumn{2}{c}{\textbf{覆盖}} & \multicolumn{2}{c}{\textbf{计划内恢复}} & \multicolumn{2}{c}{\textbf{中位/P95 (ms)}} \\
\cmidrule(lr){4-5}\cmidrule(lr){6-7}\cmidrule(lr){8-9}
 & & 任务 & B & Z & B & Z & B & Z \\
\midrule
2715 & 正常   & 6/6 & 24/24 & 24/24 & N/A(0/0) & N/A(0/0) & .223/.318 & .129/.295 \\
2716 & B断馈  & 5/6 & 19/24 & 24/24 & 1/1      & N/A(0/0) & .242/.395 & .144/.330 \\
2717 & Z断馈  & 5/6 & 24/24 & 19/24 & N/A(0/0) & 1/1      & .222/.317 & .134/.336 \\
\bottomrule
\end{tabular}
\par\smallskip
\raggedright\scriptsize
三轮均采用 \texttt{R0} 馈线直连拓扑，每轮 $300$ 槽（周期 \SI{500}{ms}：$6$ 项任务 $\times$ $K{=}3$ 与 $282$ 个等长纯时间更新槽；双侧共覆盖探针的估计差中位数为 \SIrange{0.140}{0.153}{ms}、P95 为 \SIrange{0.202}{0.264}{ms}，$n{=}24,19,19$）。按 \texttt{uc1\_time\_tail\_audit} 口径将停发后尾段 \texttt{stale}（双侧各 $1/1$）与计划内事件分离；未受干预侧无计划内失更事件，记为 \texttt{N/A (0/0)}。每轮均有 $0/6$ 项计划任务通过定时软件处理判定（外部射频参考界未校准）。表中数据为本侧 IQ 采样域留出绝对预测误差，不代表跨设备时钟同步精度。
\end{table}
